%! Linear Transformations — Unit 5, Lesson 5.3 — Warm-Up (Answer Key)
\documentclass[10pt]{article}
\usepackage{linalg-article}
\usepackage{linalg-key}   % pulls in -boxes; do NOT also load -boxes
\newcommand{\TallMath}[1]{$\displaystyle #1\rule[-1.4em]{0pt}{3.2em}$}
\newcommand{\ee}{\mathbf{e}}

\begin{document}

\pageheader{Unit 5, Lesson 5.3}{Warm-Up --- Answer Key}
\namedateperiod

\vspace{0.12in}

\begin{notesbox}{Getting Ready: Where Does a Matrix Send the Unit Square?}
\small We will run all three steps on the \emph{same} matrix
$A=\begin{bsmallmatrix} 3 & 1 \\ 0 & 2 \end{bsmallmatrix}$.

\vspace{4pt}
\textbf{1. Where do the corners go? (Lesson 1.3).} Recall $A\mathbf{x}$ is a combination of the columns,
so $A\ee_1$ and $A\ee_2$ are just the columns of $A$. Fill them in:
\[
  A\ee_1=A\begin{bmatrix} 1 \\ 0 \end{bmatrix}=\begin{bmatrix} {\color{keyred}\mathbf{3}} \\ {\color{keyred}\mathbf{0}} \end{bmatrix},
  \qquad
  A\ee_2=A\begin{bmatrix} 0 \\ 1 \end{bmatrix}=\begin{bmatrix} {\color{keyred}\mathbf{1}} \\ {\color{keyred}\mathbf{2}} \end{bmatrix}.
\]

\vspace{4pt}
\textbf{2. A $2\times2$ determinant (Lesson 5.0).} Recall
$\det\begin{bsmallmatrix} a & b\\ c & d \end{bsmallmatrix}=ad-bc$. Compute
\[
  \det A = \det\begin{bmatrix} 3 & 1 \\ 0 & 2 \end{bmatrix} = {\color{keyred}\mathbf{(3)(2)-(1)(0)=6}}.
\]

\vspace{4pt}
\textbf{3. From square to parallelogram.} The unit square (area $1$) has corners at $\mathbf{0}$,
$\ee_1$, $\ee_2$, and $\ee_1+\ee_2$. After applying $A$, those corners land on $\mathbf{0}$ and the two
columns you found in step~1 --- so the square becomes the \emph{parallelogram} with edges $(3,0)$ and
$(1,2)$. Its area turns out to equal $|\det A|$:
\[
  \text{area of the image} = {\color{keyred}\mathbf{|\det A|=6}}.
\]

\vspace{4pt}
\textbf{One-sentence prediction.} If applying $A$ turns a square of area $1$ into a shape of area
$|\det A|$, what does it do to \emph{every} area? \ansline{It multiplies every area by $|\det A|=6$ --- the determinant is the factor the transformation scales area by.}
\end{notesbox}

\begin{teachernote}
This is the whole lesson in miniature. Step~1 recalls that $A\ee_1,A\ee_2$ are the columns (Lesson~1.3),
step~2 recalls the $2\times2$ determinant (Lesson~5.0), and step~3 fuses them: the unit square lands on
the column parallelogram, whose area is $|\det A|=6$. Name the punchline aloud --- ``applying $A$
multiplies \emph{every} area by $|\det A|$'' --- because that single fact is the spine of the notes.
Watch for $ad+bc$ in step~2, and for students thinking the area is $3\cdot2=6$ ``because the box is
$3\times2$'': it is $6$, but the reason is the determinant (the $(1,2)$ column is slanted, yet the base
$\times$ height still works out to $6$).
\end{teachernote}

\end{document}
